\documentclass[11pt]{article}

% --- Series style ---
\usepackage{series}

% --- Math and layout ---
\usepackage{amsmath,amssymb,amsthm}
\usepackage{mathtools}
\usepackage{booktabs}
\usepackage{geometry}
\usepackage{hyperref}

\geometry{margin=1in}


% --- Metadata ---
\title{Superset Determinations in Modal Triplet Theory:\\Parameter Identifiability after Profile-Standard SM Closure}
\author{Peter Nero}
\date{July 2026\\Version 3}

\begin{document}
\maketitle

\begin{abstract}
We formulate the MTT superset program as a parameter-identifiability
framework. Embedded renormalized-Standard-Model equivalence closes at the
adopted one-shared-physical-primitive and measured-profile standard using
selected SMDR~v1.3 transport and a complete eight-row covariance. The
construction ledger contains one shared electroweak primitive and 15 measured
SM profile coordinates; transported couplings, nine charged Yukawa rows, the
Higgs row, covariance entries, and the 96-dimensional finite Dirac operator
are outputs and are not counted again as independent inputs. Exact source-side
structure includes the native faithful gauge group, anomaly-free chiral
carrier, completed finite triple, one-Higgs projector, and 10/10 finite
gauge/ghost spectrum rows. None of these presently replaces the measured
profile by a zero-primitive source theorem. We classify every result as a
structural theorem, calibration, profile/replay, diagnostic, or held-out
prediction and state the resulting identifiability limits.
\end{abstract}

\section*{Version 3 Revision Note}
\addcontentsline{toc}{section}{Version 3 Revision Note}
\begin{description}
\item[Supersedes] \emph{Superset Determinations in Modal Triplet Theory},
version~2.
\item[Reason] One-loop crossings, fitted zeta ratios and threshold vectors,
and the claimed $\alpha_s(M_Z)$ cross-prediction reused target-related data and
were not independent determinations. The earlier parameter discussion also
risked double-counting transported outputs as independent inputs.
\item[Resolution] Version~3 replaces the numerical claim chain by a parameter-
identifiability ledger tied to the selected SMDR profile and classifies every
quantity by theorem, calibration, replay, diagnostic, or prediction status. It
incorporates the native finite-geometry results and separates finite-replay
Yukawa closure from strict source-value emission.
\item[Retained result] The superset strategy remains useful for comparing
encodings, exposing shared primitives, and discovering common-source
constraints.
\item[Remaining boundary] Predictive promotion requires a selected continuum
source/action that replaces the measured profile and supports genuinely
held-out observables.
\end{description}

\section{Purpose of the superset program}

MTT can organize several effective encodings through shared carrier, projector,
bundle, overlap, and response data. The superset strategy is useful when it
asks which observables are jointly representable and which parameters are
identifiable from a declared input set. It is not, by itself, a method for
deriving measured constants.

We use five disjoint status classes:
\begin{enumerate}
\item \textbf{Structural theorem:} follows from selected mathematical data
without observed value selection.
\item \textbf{Calibration:} parameters are chosen to reproduce declared data.
\item \textbf{Profile/replay:} measured coordinates are transported through a
specified scheme or embedded action.
\item \textbf{Diagnostic:} a comparison tests consistency but was involved in
selection or fitting.
\item \textbf{Held-out prediction:} the compared observable was absent from
source selection, calibration, branch choice, and uncertainty construction.
\end{enumerate}

\section{Retirement of the legacy Tier-3 execution}

The former scales $\Lambda_{12}$ and $\Lambda_{23}$ were obtained from a
one-loop crossing exercise. The reported values near $4.2$--$5~\mathrm{TeV}$
were then used to extract $\zeta$ ratios, fit a common scale, choose a
minimum-norm threshold, and reconstruct $\alpha_s$. This chain is retired.

Specifically:
\begin{itemize}
\item a gauge crossing is not an MTT coherence or cutoff scale;
\item the old $0.560$ and $0.229$ ratios are not current profile values;
\item $K$ was calibrated with gravitational/electroweak information and was
not an independent prediction;
\item the minimum-norm threshold was a regularized fit selected by an external
norm, not a source theorem; and
\item $\alpha_s$ was not held out from the latent extraction and threshold
choices, so it was not a cross-prediction.
\end{itemize}

\section{Current common-scheme profile}

The selected numerical transport uses SMDR~v1.3 at
\[
Q=M_t=172.5590883453979~\mathrm{GeV}.
\]
Its gauge entries are
\[
\begin{aligned}
g_Y&=0.3585945042\pm0.0000307251,\\
g_2&=0.6475986708\pm0.0000287665,\\
g_3&=1.163427409\pm0.004036156.
\end{aligned}
\]
Fifteen locked source coordinates are transported into eight full-SM
$\overline{\mathrm{MS}}$ rows. The differentiated map emits a positive-definite
$8\times8$ covariance with all 36 symmetric entries and all 15 BCT--WZH
cross-block entries. These rows are profile transport, not MTT-derived
constants.

\section{Adopted closure and parameter count}

The final global audit closes 12/12 obligations for embedded renormalized-SM
equivalence at the adopted standard. The independent-input ledger is:
\begin{center}
\begin{tabular}{p{0.43\linewidth}p{0.12\linewidth}p{0.34\linewidth}}
\toprule
Class & Count & Interpretation \\
\midrule
Construction-side continuous primitive & 1 & shared
$P_{\mathrm{EW}}$, counted once \\
Measured SM profile coordinates & 15 & empirical common-scheme source
profile with declared covariance policy \\
Neutral extension coordinates & 2 & measured oscillation splittings, outside
the 12/12 SM baseline \\
Higgs-specific construction parameters & 0 & no separately fitted Higgs
primitive \\
\bottomrule
\end{tabular}
\end{center}

The eight SMDR output rows, nine charged Yukawa magnitude rows, finite
$96\times96$ Dirac matrix, threshold/mass-scheme rows, and propagated
covariance are transported or reconstructed outputs. They are not additional
independent inputs and must not be counted a second time. Conversely, calling
them outputs does not make the 15 measured source coordinates predictions.
The current ledger therefore makes no strict parameter-reduction claim
relative to the Standard Model.

The distinction between two closure standards matters here. The nine charged
Yukawa rows close at the accepted finite-replay/profile standard, and the
Higgs lane closes with zero Higgs-specific parameter once
$P_{\mathrm{EW}}$ is declared. At the stronger zero-primitive tier, however,
the observed magnitudes and the value of $P_{\mathrm{EW}}$ are not emitted by
selected source geometry. Thus strict source replacement remains open without
reopening the profile-standard Yukawa or Higgs result.

\section{Selected structural outputs}

The current superset contains genuine source-side structure independent of the
profile values, including:
\begin{itemize}
\item the locked 27-by-27 qutrit--Weyl/minimal matrix architecture;
\item the selected rank-two literal Cech--HYM witness, closed 2/2;
\item the internal $K_{\mathrm{threshold}}$ response ledger, closed 10/10;
\item the selected retarded $q79/F/m1$ orientation representative; and
\item the CKM $\Pi$ rows, closed at the prediction-with-uncertainty-profile
standard.
\item the native faithful
$(SU(3)\times SU(2)\times U(1))/\mathbb Z_6$ gauge group and anomaly-free
48-state chiral carrier;
\item the completed 96-dimensional finite triple at profile tier and its
single anomaly-free shared hypercharge circle;
\item the exact selected one-Higgs projector; and
\item all ten finite gauge/ghost spectrum-source rows, with zero continuous
parameters.
\end{itemize}
These achievements constrain how the measured profile can be embedded. They do
not automatically emit every profile magnitude. In particular, the normalized
finite gauge spectra are universal and therefore cannot determine the observed
relative gauge couplings without a selected continuum response.

\section{Strict-upgrade ledger}

Beyond the adopted baseline, nine stronger obligations are tracked:
\begin{center}
\begin{tabular}{cll}
\toprule
ID & Upgrade & Status \\
\midrule
U1 & zero-primitive empirical source & open \\
U2 & literal global Cech--HYM/QaSU3 & closed 2/2 \\
U3 & official joint input likelihood & partial \\
U4 & CKM prediction beyond source rows & closed at profile criterion \\
U5 & absolute neutrino mass and ontology & partial \\
U6 & strong-CP selection & partial \\
U7 & MTT-derived quantization & partial \\
U8 & constructive nonperturbative 4D QFT & partial \\
U9 & unique observed-branch selection & partial \\
\bottomrule
\end{tabular}
\end{center}
Closing U2 or U4 does not imply that the other seven upgrades close.

\section{Identifiability theorem}

\begin{theorem}[Current-ledger identifiability]
Given the selected branch, the 15-coordinate common-scheme profile, one shared
electroweak primitive, the embedded SM action, and the selected finite operator
data, the accepted observable rows are jointly reproducible with their
declared covariance. If the measured profile is removed while the current
source maps are held fixed, the physical numerical rows are not identifiable:
the exact finite data do not supply a replacement map to the missing
dimensionful and relative-coupling coordinates.
\end{theorem}

\begin{proof}
The selected transport and embedding provide the forward map and covariance at
the profile standard. The finite-replay Yukawa and Higgs constructions consume
that profile and the declared shared primitive; the native finite gauge
spectra are universal after lane normalization. No accepted current map sends
the selected continuum geometry and upper action to all 15 measured source
coordinates. Therefore the profile is sufficient for the adopted forward
execution and presently necessary for its physical numerical values. The
statement is about the current source inventory, not an impossibility theorem
for future MTT constructions.
\end{proof}

\section{Interpretation of the shared structure}

The superset does achieve something more specific than a relabeling of
Standard Model parameters. One discrete carrier architecture supports gauge,
fermion, Higgs, threshold, and flavor constructions with common projectors and
provenance. This creates cross-sector consistency tests: a proposed source
cannot be changed independently in one sector if the same object acts in the
others. The native gauge quotient, anomaly table, finite triple, and one-Higgs
projector are examples of such shared restrictions.

What the superset has not yet achieved is a generative numerical law for the
measured profile. A compact representation with shared objects may be
structurally economical without using fewer empirical coordinates at the
strict prediction tier. The appropriate comparison with the Standard Model is
therefore two-column: derived structural constraints on one side, and
empirical continuous inputs on the other. Combining those columns into a
single ``parameter count'' would obscure rather than demonstrate progress.

\section{Rules for future superset calculations}

Every future execution must publish raw inputs, units, scheme and scale,
selection rules, fitted parameters, covariance, branch choices, code, and unit
tests. A claimed held-out prediction must include a data-separation certificate
showing that the observable was absent from all source and calibration paths.
No matching point, internal spectral gap, or gauge crossing may be promoted to
a physical cutoff without a separate source theorem.

\section{Conclusion}

The superset program survives, but in a sharper form. Its current success is
the coherent, reproducible embedding of the renormalized SM and the reduction
of many sector-specific constructions to shared finite operator data. The
native gauge, chiral, finite-triple, one-Higgs, and finite-spectrum results add
real source-side restrictions without changing the empirical-input count. The
old crossing-scale calibration and its downstream numerical claims are
retired.

The next frontier is not another backfit. It is source promotion: emit the ten
magnitude rows, relative gauge metrics, electroweak normalization, and
physical threshold/mass values from one selected continuum action, or retain
their measured source coordinates transparently. That distinction is the
correct measure of progress toward strict no-knob closure.

% BEGIN MTT MANAGED COMPUTATIONAL EVIDENCE
\section*{Computational Evidence and Reproducibility}
The listed packets are used directly as the finite family of profile, exact, and certified-numerical examples on which the identifiability analysis operates. Their values are not rederived here, and a profile row remains profile replay. The open ledger states the unresolved reduction from those effective coordinates to a strict no-knob source.

The referenced rows are frozen to the curated results repository state identified below. Hashes are grouped in eight-character blocks for line breaking.
\begin{quote}\small
\textbf{Repository:} \url{https://github.com/PeterNero/mtt-results-repro}\\
\textbf{Commit:} \texttt{31247ebb 5c22f3fb b5443024 365433c6 ee0bff4a}\\
\textbf{Manifest:} \path{release/result_manifest.json}\\
\textbf{Manifest SHA-256:}\\
\texttt{fb399689 60b00584 631dbf53 1a708e18}\\
\texttt{ef928d6b 6d935119 c185d7f6 32b1e7cd}
\end{quote}

Tier labels are quoted verbatim from that manifest. A row used directly supports only the specific computational statement identified above; a corpus-state cross-check does not prove this paper's local theorems; and an open row is evidence of an unresolved obligation, never of closure.

\par\begingroup\dimen0=\pagegoal\advance\dimen0 by -\pagetotal\ifdim\dimen0<7\baselineskip\newpage\fi\endgroup
\paragraph{Rows used directly in this paper.}
\begin{itemize}
\setlength{\itemsep}{0.25em}
\item \path{A01/charged_yukawa_higgs_profile} (\emph{profile replay}).\par Versioned Yu, Yd, Ye and lambda\_H profile packet.
\item \path{A14/ckm_prediction_profile} (\emph{numeric certified}).\par Three selected CKM profile rows and uncertainty comparison.
\item \path{A01/current_global_lock} (\emph{profile replay}).\par Current non-looping global status and source-certificate map.
\item \path{A01/direct_k_higgs_row} (\emph{derived exact}).\par Promoted direct K\_threshold.Omega\_H.lambda row.
\item \path{A22/e6_qpsi_qcd_anomaly} (\emph{derived exact}).\par E6 Qpsi matter/exotic QCD anomaly cancellation audit.
\item \path{A04/final_12_of_12_audit} (\emph{profile replay}).\par Twelve-obligation embedded renormalized-SM equivalence audit.
\item \path{A19/hym_wiener_contraction} (\emph{numeric certified}).\par Weighted-theta Fourier-tail and Wiener contraction certificate.
\item \path{A40/neutral_two_primitive_profile} (\emph{profile replay}).\par Measured two-splitting neutrino profile, masses and Dirac Yukawa rows.
\item \path{A02/precision_15_source_transport} (\emph{profile replay}).\par Fifteen measured source coordinates, Jacobian and covariance transport.
\item \path{A02/precision_8x8_workspace} (\emph{profile replay}).\par Eight-coordinate SMDR output with positive-definite 8x8 covariance.
\item \path{A01/qutrit_weyl_27_matrix} (\emph{derived exact}).\par Sparse 27x27 qutrit-Weyl left-action realization.
\item \path{A01/strict_pew_row} (\emph{derived exact}).\par Promoted P\_EW source row at the declared one-shared-primitive standard.
\end{itemize}

\par\begingroup\dimen0=\pagegoal\advance\dimen0 by -\pagetotal\ifdim\dimen0<7\baselineskip\newpage\fi\endgroup
\paragraph{Open boundary (not evidence of closure).}
\begin{itemize}
\setlength{\itemsep}{0.25em}
\item \path{A05/strict_upgrade_ledger} (\emph{open}).\par Current 2/9 strict no-knob upgrade ledger.
\end{itemize}

No imported row changes theorem ownership or promotes a neighboring claim: all local statements retain their stated hypotheses, domains, and limitations.
% END MTT MANAGED COMPUTATIONAL EVIDENCE

\begin{thebibliography}{99}

\bibitem{MTT-Admissibility}
P.~Nero,
\emph{Modal Triplet Theory: Admissibility, Encodings, and the Structure of Physical Description},
Zenodo preprint, January 2026.
\url{https://doi.org/10.5281/zenodo.18255621}

\bibitem{MTT-Foundation}
P.~Nero,
\emph{Modal Triplet Theory: Foundation},
Zenodo preprint, September 2025.
\url{https://doi.org/10.5281/zenodo.16949762}

\bibitem{MTT-FixedPoints}
P.~Nero,
\emph{Fixed Points I--VI: Complete Coherence Spine},
Zenodo preprints, August 2025.
\url{https://doi.org/10.5281/zenodo.16948748}

\bibitem{MTT-ProjectionAdmissibility}
P.~Nero,
\emph{The Projection--Admissibility Principle: Structural Constraints on Effective Physical Description},
Zenodo preprint, January 2026.
\url{https://doi.org/10.5281/zenodo.18255838}

\bibitem{MTT-Closure}
P.~Nero,
\emph{Closure and Inevitability in Modal Triplet Theory},
Zenodo preprint, January 2026.
\url{https://doi.org/10.5281/zenodo.18255510}

\bibitem{MTT-CoherenceCapacity}
P.~Nero,
\emph{Coherence Capacity as the Fundamental Resource of Effective Physics},
Zenodo preprint, January 2026.
\url{https://doi.org/10.5281/zenodo.18255905}

\bibitem{MTT-CoherenceDynamics}
P.~Nero,
\emph{Dynamics of Coherence Capacity: Transport, Concentration, and Exhaustion},
Zenodo preprint, January 2026.
\url{https://doi.org/10.5281/zenodo.18256048}

\bibitem{MTT-QM}
P.~Nero,
\emph{Modal Triplet Theory: From MTT to Quantum Mechanics},
Zenodo preprint, September 2025.
\url{https://doi.org/10.5281/zenodo.17074246}

\bibitem{MTT-QFT}
P.~Nero,
\emph{From MTT to Quantum Field Theory},
Zenodo preprint, 2025.
\url{https://doi.org/10.5281/zenodo.17068816}

\bibitem{MTT-GR}
P.~Nero,
\emph{Modal Triplet Theory: From MTT to General Relativity},
Zenodo preprint, October 2025.
\url{https://doi.org/10.5281/zenodo.16950597}

\bibitem{MTT-UVQG}
P.~Nero,
\emph{Modal Triplet Theory: From MTT to a UV-Finite, Unitary Quantum Gravity},
Zenodo preprint, 2025.
\url{https://doi.org/10.5281/zenodo.17077671}

\bibitem{MTT-Measurement}
P.~Nero,
\emph{Measurement as Disturbance and Stabilization in Modal Triplet Theory},
Zenodo preprint, 2025.
\url{https://doi.org/10.5281/zenodo.17177404}

\bibitem{MTT-Irreversibility}
P.~Nero,
\emph{Projection, Probability, and Irreversibility: Shadow Bridges Between Measurement, Black Holes, and Cosmology in Modal Triplet Theory},
Zenodo preprint, January 2026.
\url{https://doi.org/10.5281/zenodo.18256408}

\bibitem{MTT-Bell-Beables}
P.~Nero,
\emph{Modal Fixed Points, Bell's Beables, and the Limits of Factorization},
Zenodo preprint, 2025.
\url{https://doi.org/10.5281/zenodo.17076300}

\bibitem{MTT-Bell-Temporal}
P.~Nero,
\emph{Temporal Bell Inequalities and Global Consistency in Modal Triplet Theory},
Zenodo preprint, August 2025.
\url{https://doi.org/10.5281/zenodo.18208884}

\bibitem{MTT-Stochastic}
P.~Nero,
\emph{From Modal Triplet Theory to Indivisible Stochastic Processes: A First-Principles, Fully Rigorous Derivation},
Zenodo preprint, January 2026.
\url{https://doi.org/10.5281/zenodo.18254862}

\end{thebibliography}
\end{document}
